\documentclass[a4paper,12pt]{article}
\usepackage[T2A]{fontenc}
\usepackage[utf8]{inputenc}
\usepackage[russian]{babel}
\usepackage{PTSans}
\renewcommand{\familydefault}{\sfdefault}
\usepackage{amsmath,amssymb,mathtools}
\usepackage{icomma}
\usepackage[a4paper,left=21mm,right=21mm,top=19mm,bottom=19mm,headheight=26pt,headsep=4mm]{geometry}
\usepackage{microtype}
\usepackage{tikz}
\usetikzlibrary{arrows.meta,positioning,shapes.geometric,calc}
\usepackage{booktabs,tabularx}
\usepackage{enumitem}
\usepackage{fancyhdr}
\usepackage{setspace}
\usepackage{xcolor}
\usepackage{hyperref}
\definecolor{accent}{HTML}{087C82}
\definecolor{darkaccent}{HTML}{174F59}
\definecolor{textgray}{HTML}{333D43}
\definecolor{softgray}{HTML}{61717A}
\definecolor{lightline}{HTML}{D3DFE1}
\definecolor{warmgray}{HTML}{F2F5F5}
\hypersetup{colorlinks=true,linkcolor=darkaccent,urlcolor=darkaccent,
 pdfauthor={Лёвин Артём Александрович},
 pdftitle={Баланс сил. Второй закон Ньютона. 21.10.26},
 pdfsubject={Физика, подготовка к ОГЭ, чек-лист для Володи Хачатуряна}}
\setstretch{1.14}
\setlength{\parindent}{0pt}
\setlength{\parskip}{3pt}
\setlist[itemize]{leftmargin=*,topsep=3pt,itemsep=3pt,parsep=0pt}
\setlist[enumerate]{leftmargin=*,topsep=5pt,itemsep=11pt,parsep=0pt}
\pagestyle{fancy}
\fancyhf{}
\fancyhead[L]{\small\color{softgray}ФИЗИКА \(\cdot\) ОГЭ}
\fancyhead[R]{\small\color{softgray}21.10.26}
\fancyfoot[C]{\small\color{softgray}\thepage}
\renewcommand{\headrulewidth}{0.35pt}
\renewcommand{\headrule}{\hbox to\headwidth{\color{lightline}\leaders\hrule height\headrulewidth\hfill}}
\newcommand{\pagetitle}[1]{{\LARGE\bfseries\color{darkaccent}#1}\par\vspace{3pt}\noindent{\color{lightline}\rule{\linewidth}{0.55pt}}\par\vspace{7pt}}
\newcommand{\heading}[1]{\vspace{7pt}{\large\bfseries\color{darkaccent}#1}\par\vspace{2pt}}
\newcommand{\smallhead}[1]{\vspace{4pt}{\bfseries\color{darkaccent}#1}\par\vspace{2pt}}
\newcommand{\focus}[1]{\textcolor{darkaccent}{\textbf{#1}}}
\newcommand{\unit}[1]{\,\text{#1}}
\newcommand{\checkitem}{\item[$\square$]}
\newcommand{\fres}{F_{\text{сопр}}}
\newcommand{\farch}{F_{\text{A}}}
\newcommand{\vproj}{a_y}
\tikzset{
 bodybox/.style={rectangle,draw=darkaccent,fill=warmgray,rounded corners=1.5pt,minimum width=17mm,minimum height=12mm,line width=0.9pt},
 force/.style={-{Stealth[length=3mm]},line width=1.5pt,draw=accent},
 gravity/.style={-{Stealth[length=3mm]},line width=1.5pt,draw=darkaccent},
 accel/.style={-{Stealth[length=2.8mm]},line width=1.4pt,draw=black,dashed},
 ax/.style={-{Stealth[length=2.5mm]},line width=1pt,draw=softgray},
 flowstart/.style={ellipse,draw=darkaccent,line width=.9pt,align=center,minimum width=35mm,minimum height=9mm,fill=warmgray},
 flowaction/.style={rounded corners=2pt,rectangle,draw=darkaccent,line width=.9pt,align=center,text width=99mm,minimum height=12mm,inner sep=4pt},
 flowdecision/.style={diamond,aspect=2.45,draw=darkaccent,line width=.9pt,align=center,text width=47mm,inner sep=2pt},
 flowarr/.style={-{Stealth[length=2.8mm]},line width=1pt,draw=textgray}
}
\begin{document}
\thispagestyle{empty}
\begin{center}
\vspace*{35mm}
{\small\bfseries\color{accent}ФИЗИКА \(\cdot\) ПОДГОТОВКА К ОГЭ}\par
\vspace{17mm}
{\Huge\bfseries\color{darkaccent}Чек-лист}\par
\vspace{9mm}
{\LARGE\bfseries Баланс сил}\par
\vspace{3mm}
{\Large Второй закон Ньютона в задачах}\par
\vspace{10mm}
{\normalsize\color{textgray}Лифт, подвешенный груз, сопротивление воздуха,}\par
{\normalsize\color{textgray}сила Архимеда и связь с кинематикой}\par
\vspace{18mm}
{\large 21 октября 2026 года}\par
\vfill
{\normalsize\color{textgray}Лёвин Артём Александрович}\par
{\normalsize\color{textgray}эксклюзивно для Володи Хачатуряна}\par
\vspace{22mm}
\begin{tikzpicture}[x=1cm,y=1cm]
\draw[accent!70,line width=1.2pt] (-5,0) .. controls (-2.7,0.85) and (-1.3,-0.85) .. (1,0) .. controls (2.8,0.65) and (4.1,0.45) .. (5,0.1);
\end{tikzpicture}
\end{center}
\clearpage

\pagetitle{1. Главная идея: силы и проекции}
\begin{itemize}
\checkitem \focus{Второй закон Ньютона} в инерциальной системе отсчёта:
\[\boxed{\sum\vec F=m\vec a},\qquad\boxed{\sum F_y=ma_y}.\]
В левой части --- \emph{все силы, действующие на выбранное тело}; в правой --- масса, умноженная на \emph{проекцию ускорения}.
\checkitem \focus{Знак проекции} определяется направлением вектора относительно выбранной оси. Совпадает с осью --- «\(+\)»; противоположен --- «\(-\)». Знак \emph{не выбирают} по тому, какая сила больше.
\checkitem \focus{Скорость и ускорение --- разные векторы.} Тело может двигаться вверх, но ускоряться вниз (например, при торможении подъёма). Ускорение определяется равнодействующей сил.
\checkitem При \(a_y=0\) сумма проекций сил равна нулю. Тело может \emph{покоиться} или двигаться \emph{равномерно и прямолинейно}.
\end{itemize}
\heading{Какие силы изображаем на рисунке}
\renewcommand{\arraystretch}{1.26}
\begin{tabularx}{\linewidth}{@{}>{\bfseries}p{35mm}X@{}}
\toprule
Сила & Куда направлена и когда появляется \\
\midrule
Тяжести \(mg\) & Вертикально вниз; \(g\approx10\unit{м/с\textsuperscript{2}}\) в задачах этого занятия. \\
Реакция опоры \(N\) & Перпендикулярно опоре, от опоры к телу; контактная сила. \\
Натяжение нити \(T\) & Вдоль натянутой нити, от тела к точке крепления. \\
Сопротивление \(\fres\) & Против скорости тела относительно среды, если сопротивление учитывается. \\
Архимедова \(\farch\) & Вертикально вверх в покоящейся среде: \(\farch=\rho_{\text{ср}}gV_{\text{погр}}\). \\
\bottomrule
\end{tabularx}
\heading{Мини-проверка знаков}
Выбрали ось \(Oy\) \emph{вверх}; \(N\) вверх, \(mg\) вниз, ускорение вниз:
\[\boxed{N-mg=-ma}.\]
Теперь направим \(Oy\) \emph{вниз}, не меняя физических условий:
\[\boxed{mg-N=ma}.\]
\focus{Оба уравнения описывают один и тот же случай.} В обоих \(N=m(g-a)\) при \(0\le a\le g\).
\smallhead{Привычка перед каждой задачей}
Изобразите силы \(\rightarrow\) выберите ось \(\rightarrow\) расставьте знаки \(\rightarrow\) запишите \(\sum F_y=ma_y\).
\clearpage

\pagetitle{2. Тело на полу движущегося лифта}
\begin{center}
\begin{tikzpicture}[scale=0.95]
\draw[thick,softgray] (-2.3,-1.3) -- (-2.3,2.3) -- (2.3,2.3) -- (2.3,-1.3);
\draw[line width=1.3pt] (-2.25,-1.3)--(2.25,-1.3);
\node[bodybox] (b) at (0,-0.63) {$m$};
\draw[force] (0,-.02)--(0,1.22) node[above] {$N$};
\draw[gravity] (.42,-.72)--(.42,-2.05) node[below] {$mg$};
\draw[accel] (-1.4,-.32)--(-1.4,1.35) node[above] {$\vec a$};
\draw[ax] (2.95,-1.0)--(2.95,1.08) node[above] {$y$};
\node[font=\small,text=softgray] at (-1.65,-2.0) {лифт};
\end{tikzpicture}
\end{center}
\begin{itemize}
\checkitem На тело действуют \emph{две силы}: вниз \(mg\) и вверх \(N\). Ускорение \(\vec a\) --- результат действия сил, его не добавляют в сумму сил.
\checkitem Если ускорение направлено \focus{вверх}, при оси вверх:
\[N-mg=ma\quad\Longrightarrow\quad\boxed{N=m(g+a)}.\]
\checkitem Если ускорение направлено \focus{вниз} и ось вверх:
\[N-mg=-ma\quad\Longrightarrow\quad\boxed{N=m(g-a)}.\]
При сохранении контакта с полом \(0\le a\le g\); при \(a=g\) реакция опоры обращается в нуль. Если пол ускоряется вниз быстрее свободного падения, тело отрывается от него и формула контакта уже неприменима.
\checkitem При \(a=0\): \(N=mg\) независимо от того, покоится ли лифт или движется с постоянной скоростью.
\end{itemize}
\smallhead{Пример: лифт ускоряется вверх}
Груз \(m=3\unit{кг}\), ускорение лифта \(a=2\unit{м/с\textsuperscript{2}}\). Тогда
\[N=3(10+2)=\boxed{36\unit{Н}}.\]
\smallhead{Пример: лифт ускоряется вниз}
Тот же груз и то же по модулю ускорение:
\[N=3(10-2)=\boxed{24\unit{Н}}.\]
\focus{Важно:} направление \emph{движения} лифта ещё не сообщает направление ускорения. Лифт может подниматься и одновременно замедляться --- ускорение в таком случае вниз.
\clearpage

\pagetitle{3. Груз на нити: сила натяжения}
\begin{center}
\begin{tikzpicture}[scale=.96]
\draw[line width=1.2pt,softgray] (-2.6,1.85)--(2.6,1.85);
\foreach \x in {-2.45,-2.05,...,2.4} {\draw[softgray] (\x,1.85)--(\x+.22,2.07);}
\draw[line width=1pt,textgray] (0,1.85)--(0,0);
\node[bodybox] (b) at (0,-.35) {$m$};
\draw[force] (0,-.14)--(0,1.13) node[right] {$T$};
\draw[gravity] (.42,-.54)--(.42,-1.82) node[right] {$mg$};
\draw[accel] (-1.55,-.45)--(-1.55,.86) node[above] {$\vec a$};
\draw[ax] (2,-1.05)--(2,.85) node[above] {$y$};
\end{tikzpicture}
\end{center}
\begin{itemize}
\checkitem Для вертикальной \emph{натянутой} нити с грузом: \(T\) вверх, \(mg\) вниз. Ось удобно направить вверх.
\checkitem Груз в покое или движется равномерно:
\[T-mg=0\quad\Longrightarrow\quad\boxed{T=mg}.\]
\checkitem Ускорение груза направлено вверх:
\[T-mg=ma\quad\Longrightarrow\quad\boxed{T=m(g+a)}.\]
\checkitem Ускорение направлено вниз:
\[T-mg=-ma\quad\Longrightarrow\quad\boxed{T=m(g-a)}.\]
В этом случае натянутая нить возможна при \(a<g\). При \(a=g\) получаем \(T=0\) --- груз в свободном падении, натяжение отсутствует. Нить не может \emph{толкать} груз: отрицательное \(T\) физически недопустимо.
\end{itemize}
\smallhead{Пример занятия: груз массой \(2\unit{кг}\)}
Если система ускоряется \emph{вверх} с \(a=20\unit{м/с\textsuperscript{2}}\), то
\[T=2(10+20)=\boxed{60\unit{Н}}.\]
При \emph{таком же ускорении вниз} простой модели висящего на натянутой нити груза уже недостаточно: формула дала бы \(T=-20\unit{Н}\). Это указывает на потерю натяжения; без иных сил груз не может ускоряться вниз быстрее \(g\).
\heading{Что объединяет лифт и подвешенный груз}
Если тело поддерживают \emph{снизу}, ищем \(N\); если его держит \emph{нить сверху}, ищем \(T\). В обоих случаях решение строится из двух сил и одного уравнения Ньютона.
\clearpage

\pagetitle{4. Падение и сопротивление воздуха}
\begin{center}
\begin{tikzpicture}[scale=.98]
\node[bodybox] at (0,0) {$m$};
\draw[gravity] (-.22,-.45)--(-.22,-2.0) node[left] {$mg$};
\draw[force] (.28,.43)--(.28,1.85) node[right] {$\fres$};
\draw[ax] (2,1.25)--(2,-1.1) node[below] {$y$};
\draw[accel] (-1.7,.3)--(-1.7,-1.15) node[below] {$\vec a$};
\node[text=softgray,font=\small] at (-2.25,.8) {падение вниз};
\end{tikzpicture}
\end{center}
\begin{itemize}
\checkitem Если тело падает вниз, сопротивление воздуха (при его учёте) направлено \focus{вверх}. Сила тяжести \(mg\) всегда направлена вертикально вниз.
\checkitem Выбирая ось вниз, записываем общий баланс:
\[\boxed{mg-\fres=ma_y}.\]
Здесь \(a_y\) --- \emph{алгебраическая} проекция ускорения; она может быть положительной, нулевой или отрицательной.
\checkitem Если падение \emph{равномерное}: \(a_y=0\), поэтому
\[mg-\fres=0\quad\Longrightarrow\quad\boxed{\fres=mg}.\]
Это возможно при ненулевой скорости: отсутствие ускорения \emph{не означает} отсутствие движения.
\checkitem Если тело \emph{ускоряется вниз}:
\[\boxed{a=\frac{mg-\fres}{m}=g-\frac{\fres}{m}},\qquad mg>\fres.\]
Если \(mg<\fres\) при скорости вниз, ускорение направлено вверх: падающее тело \emph{замедляется}.
\checkitem \focus{Модельный договор:} в расчётных задачах этой страницы считаем сопротивление постоянным, если дан его модуль; выталкивающей силой воздуха пренебрегаем, если об этом сказано. В реальности сопротивление обычно зависит от скорости.
\end{itemize}
\smallhead{Пример: равномерное падение}
Парашютист с суммарной массой \(70\unit{кг}\) опускается с постоянной скоростью. При \(g=10\unit{м/с\textsuperscript{2}}\):
\[\fres=mg=70\cdot10=\boxed{700\unit{Н}}.\]
\smallhead{Проверка здравым смыслом}
\(mg>\fres\) --- ускорение вниз; \(mg=\fres\) --- нулевое ускорение; \(mg<\fres\) --- ускорение вверх. Направление текущей скорости во всех трёх случаях может оставаться вниз.
\clearpage

\pagetitle{5. Воздухоплавание: три силы}
\begin{itemize}
\checkitem Помимо тяжести \(mg\) и сопротивления \(\fres\), учитываем \focus{силу Архимеда}:
\[\boxed{\farch=\rho_{\text{ср}}gV_{\text{погр}}}.\]
\(\rho_{\text{ср}}\) --- плотность воздуха или жидкости (\(\unit{кг/м\textsuperscript{3}}\)); \(V_{\text{погр}}\) --- объём вытесненной среды, равный объёму погружённой части тела (\(\unit{м\textsuperscript{3}}\)).
\checkitem \(\farch\) направлена \emph{вверх}; \(mg\) --- \emph{вниз}; \(\fres\) --- \emph{против скорости} относительно среды.
\end{itemize}
\begin{center}
\begin{tikzpicture}[scale=.9]
\node[font=\bfseries] at (-3,1.9) {Подъём с постоянной скоростью};
\node[bodybox] at (-3,0) {$m$};
\draw[force] (-3.15,.4)--(-3.15,1.5) node[left] {$\farch$};
\draw[gravity] (-3.35,-.4)--(-3.35,-1.65) node[left] {$mg$};
\draw[gravity] (-2.68,-.4)--(-2.68,-1.32) node[right] {$\fres$};
\node[font=\small,text=softgray] at (-3,-2.05) {$v\uparrow,\ a=0$};
\node[font=\bfseries] at (3,1.9) {Спуск с постоянной скоростью};
\node[bodybox] at (3,0) {$m$};
\draw[force] (2.65,.4)--(2.65,1.35) node[left] {$\farch$};
\draw[force] (3.35,.4)--(3.35,1.25) node[right] {$\fres$};
\draw[gravity] (3,-.4)--(3,-1.7) node[below] {$mg$};
\node[font=\small,text=softgray] at (3,-2.05) {$v\downarrow,\ a=0$};
\end{tikzpicture}
\end{center}
\heading{Случай A. Тело равномерно поднимается}
Направим ось вверх. Сопротивление направлено вниз:
\[\farch-mg-\fres=0\quad\Longrightarrow\quad\boxed{\farch=mg+\fres}.\]
\heading{Случай Б. Тело равномерно опускается}
Направим ось вверх. Сопротивление направлено вверх:
\[\farch+\fres-mg=0\quad\Longrightarrow\quad\boxed{mg=\farch+\fres}.\]
\heading{Если движение ускоренное}
Правая часть \emph{не обязана} равняться нулю. При оси вверх:
\[\boxed{\farch-mg\mp\fres=ma_y},\]
где знак «\(-\)» берут при скорости вверх, «\(+\)» --- при скорости вниз. Знак \(a_y\) устанавливают отдельно по направлению \emph{ускорения}.
\focus{Не путайте:} знак сопротивления определяется направлением скорости, а знак \(ma_y\) --- направлением ускорения.
\clearpage

\pagetitle{6. Соединяем динамику и кинематику}
\heading{Типовая задача, разобранная на занятии}
Тело массой \(m=0{,}1\unit{кг}\) отпускают без начальной скорости с высоты \(h=20\unit{м}\). Во время падения вверх действует постоянная сила сопротивления воздуха \(\fres=0{,}2\unit{Н}\). Принять \(g=10\unit{м/с\textsuperscript{2}}\), силой Архимеда пренебречь. \focus{Найдите время падения.}
\smallhead{Шаг 1. Намечаем связь с кинематикой}
Поскольку тело падает без изменения направления с постоянным ускорением, при оси вниз:
\[h=v_0t+\frac{at^2}{2},\qquad v_0=0,\qquad\boxed{h=\frac{at^2}{2}}.\]
Чтобы найти \(t\), сначала определим \(a\).
\smallhead{Шаг 2. Находим ускорение из второго закона Ньютона}
Ось направим вниз. Сила тяжести положительна, сопротивление отрицательно:
\[mg-\fres=ma.\]
\[0{,}1\cdot10-0{,}2=0{,}1a
\quad\Longrightarrow\quad0{,}8=0{,}1a
\quad\Longrightarrow\quad\boxed{a=8\unit{м/с\textsuperscript{2}}}.\]
\smallhead{Шаг 3. Вычисляем время}
\[20=\frac{8t^2}{2}=4t^2
\quad\Longrightarrow\quad t^2=5
\quad\Longrightarrow\quad\boxed{t=\sqrt5\unit{с}\approx2{,}24\unit{с}}.\]
Корень \(t=-\sqrt5\) отбрасываем: физическое время после начала движения неотрицательно.
\smallhead{Шаг 4. Контролируем результат}
\begin{itemize}
\checkitem \(mg=1\unit{Н}>0{,}2\unit{Н}\), поэтому ускорение действительно направлено вниз.
\checkitem \(a=8\unit{м/с\textsuperscript{2}}<g\): сопротивление уменьшает ускорение падения.
\checkitem \(h=at^2/2=8\cdot5/2=20\unit{м}\): высота восстановилась.
\end{itemize}
\heading{Обобщение при \(v_0=0\)}
Если \(mg>\fres\), \(\fres\) постоянно, \(h>0\), то
\[\boxed{a=g-\frac{\fres}{m}},\qquad
\boxed{t=\sqrt{\frac{2h}{g-\fres/m}}}.\]
\focus{Важное ограничение:} формула описывает падение вниз из состояния покоя только при положительном ускорении вниз, постоянных силах и отсутствии иных существенных воздействий.
\clearpage

\pagetitle{7. Алгоритм: как решать задачу на баланс сил}
\vspace{12pt}
\begin{center}
\begin{tikzpicture}[node distance=10.4mm]
\node[flowstart] (start) {Начало задачи};
\node[flowaction,below=of start] (choose) {Выделить одно тело и изобразить \emph{все} силы, приложенные к нему};
\node[flowaction,below=of choose] (axis) {Выбрать положительное направление оси и указать направление ускорения};
\node[flowaction,below=of axis] (project) {Назначить знаки проекций каждой силы и ускорения относительно оси};
\node[flowaction,below=of project] (newton) {Записать \(\sum F_y=ma_y\) и найти искомую силу либо ускорение};
\node[flowdecision,below=15mm of newton] (question) {Нужно найти время или путь?};
\node[flowaction,below=16mm of question] (kin) {При постоянном ускорении применить уравнение кинематики};
\node[flowaction,below=10mm of kin] (check) {Проверить знак, единицы, допустимость результата и исходные условия};
\node[flowstart,below=9mm of check] (finish) {Ответ};
\draw[flowarr] (start)--(choose);
\draw[flowarr] (choose)--(axis);
\draw[flowarr] (axis)--(project);
\draw[flowarr] (project)--(newton);
\draw[flowarr] (newton)--(question);
\draw[flowarr] (question)--node[right,near start,font=\small]{да}(kin);
\draw[flowarr] (kin)--(check);
\draw[flowarr] (check)--(finish);
\draw[flowarr] (question.east)--++(2.85,0) node[above,pos=.35,font=\small]{нет}|-(check.east);
\end{tikzpicture}
\end{center}
\clearpage

\pagetitle{8. Частые ошибки и самостоятельная проверка}
\heading{Ошибки, которые особенно важно исключить}
\begin{itemize}
\checkitem \focus{«Большая сила получает плюс».} Неверно: знак задаёт направление оси. При смене оси все проекции последовательно меняют знак.
\checkitem \focus{«Тело движется вверх, значит, ускорение вверх».} Неверно: при подъёме тело может замедляться, тогда ускорение вниз.
\checkitem \focus{«Есть скорость, значит, справа \(ma\ne0\)».} Неверно: при постоянной скорости \(a=0\), сумма сил равна нулю.
\checkitem \focus{«Сопротивление воздуха всегда вверх».} Только при движении тела вниз относительно воздуха; при движении вверх сопротивление направлено вниз.
\checkitem \focus{«\(N\) и \(T\) всегда равны \(mg\)».} Равенство характерно для рассмотренных случаев с нулевым вертикальным ускорением и только соответствующими двумя силами.
\checkitem \focus{«Отрицательное натяжение нормально».} Неверно: нить не толкает; отрицательный расчётный результат указывает, что предположение о натяжении невозможно в выбранной модели.
\checkitem \focus{«Достаточно написать \(mg-F=ma\) без рисунка».} При изменении направления движения, оси или дополнительной силы легко получить неверные знаки.
\checkitem \focus{«Сила Архимеда зависит от плотности самого тела».} В \(\farch=\rho_{\text{ср}}gV_{\text{погр}}\) входит \emph{плотность среды}.
\end{itemize}
\heading{Контрольные вопросы}
\begin{itemize}
\checkitem Могу ли я словами объяснить знак \emph{каждого} слагаемого второго закона Ньютона?
\checkitem Могу ли я перезаписать баланс сил, направив ось в противоположную сторону?
\checkitem Могу ли я отличить нулевое ускорение от нулевой скорости?
\checkitem Могу ли я по \(\vec v\) найти направление силы сопротивления, а по \(\sum\vec F\) --- направление ускорения?
\checkitem Могу ли я сначала получить ускорение из сил, затем найти время из кинематики?
\end{itemize}
\smallhead{Краткая памятка формул}
\[
\boxed{\sum F_y=ma_y},\quad
\boxed{F_{\text{тяж}}=mg},\quad
\boxed{\farch=\rho_{\text{ср}}gV_{\text{погр}}},\quad
\boxed{s=v_0t+\frac{at^2}{2}}.
\]
\clearpage

\pagetitle{9. Тренировочный блок --- Путь А}
\textit{10 задач по нарастающей сложности. Во всех задачах примите \(g=10\unit{м/с\textsuperscript{2}}\); если не сказано иное, влиянием воздуха и другими силами пренебрегите. На каждом рисунке или мысленной схеме обозначайте силы и направление оси.}
\begin{enumerate}[label=\textbf{\arabic*.},leftmargin=8mm,itemsep=9pt]
\item На тело массой \(2\unit{кг}\) действуют сила \(26\unit{Н}\) вверх и сила \(20\unit{Н}\) вниз. Ось \(Oy\) направлена вверх. Запишите второй закон Ньютона в проекциях, определите \(a_y\) и направление ускорения.
\item Груз массой \(4\unit{кг}\) стоит на полу лифта. Лифт ускоряется \emph{вверх} с \(a=2\unit{м/с\textsuperscript{2}}\). Найдите силу реакции опоры.
\item На полу лифта находится ящик массой \(5\unit{кг}\). Лифт ускоряется \emph{вниз} с \(a=3\unit{м/с\textsuperscript{2}}\). Найдите силу реакции опоры; запишите уравнение при оси \emph{вниз}.
\item Лампа массой \(2{,}5\unit{кг}\) висит неподвижно на вертикальной нити. Найдите силу натяжения нити.
\item Груз массой \(1{,}2\unit{кг}\) подвешен на нити в кабине. Кабина ускоряется вверх с \(a=1{,}5\unit{м/с\textsuperscript{2}}\). Найдите силу натяжения.
\item Груз массой \(3\unit{кг}\) висит на вертикальной нити. Груз и точка подвеса ускоряются вниз с \(a=4\unit{м/с\textsuperscript{2}}\); нить натянута. Найдите натяжение.
\item Тело массой \(0{,}8\unit{кг}\) падает \emph{с постоянной скоростью}. Найдите модуль силы сопротивления воздуха, считая, что действуют только тяжесть и сопротивление.
\item Тело массой \(0{,}5\unit{кг}\) падает вниз. Сила сопротивления равна \(1\unit{Н}\) и направлена вверх. Найдите ускорение и укажите его направление.
\item Воздушный шар с суммарной массой \(0{,}2\unit{кг}\) поднимается \emph{равномерно} в неподвижном воздухе. Сопротивление равно \(0{,}4\unit{Н}\). При \(\rho_{\text{возд}}=1{,}25\unit{кг/м\textsuperscript{3}}\) найдите объём вытесненного воздуха.
\item Тело массой \(0{,}25\unit{кг}\) отпускают без начальной скорости с высоты \(16\unit{м}\). Постоянная сила сопротивления равна \(0{,}5\unit{Н}\) и направлена вверх. Найдите \emph{ускорение и время падения}. Силой Архимеда пренебречь.
\end{enumerate}
\clearpage

\pagetitle{10. Ответы к тренировочному блоку}
\begin{tabularx}{\linewidth}{@{}>{\bfseries}p{12mm}X@{}}
\toprule
№ & Ответ и ключевой шаг проверки \\
\midrule
1 & \(26-20=2a_y\), откуда \(a_y=3\unit{м/с\textsuperscript{2}}\); ускорение вверх.\\[5pt]
2 & \(N=m(g+a)=4(10+2)=\boxed{48\unit{Н}}.\)\\[5pt]
3 & При оси вниз: \(mg-N=ma\); \(N=5(10-3)=\boxed{35\unit{Н}}.\)\\[5pt]
4 & \(T=mg=2{,}5\cdot10=\boxed{25\unit{Н}}.\)\\[5pt]
5 & \(T=m(g+a)=1{,}2(10+1{,}5)=\boxed{13{,}8\unit{Н}}.\)\\[5pt]
6 & \(T=m(g-a)=3(10-4)=\boxed{18\unit{Н}}.\)\\[5pt]
7 & \(a=0\), поэтому \(\fres=mg=0{,}8\cdot10=\boxed{8\unit{Н}}.\)\\[5pt]
8 & \(ma=mg-\fres=5-1=4\unit{Н}\); \(a=\boxed{8\unit{м/с\textsuperscript{2}}}\) вниз.\\[5pt]
9 & \(\farch=mg+\fres=2+0{,}4=2{,}4\unit{Н}\); \(V=2{,}4/(1{,}25\cdot10)=\boxed{0{,}192\unit{м\textsuperscript{3}}}.\)\\[5pt]
10 & \(a=(mg-\fres)/m=(2{,}5-0{,}5)/0{,}25=\boxed{8\unit{м/с\textsuperscript{2}}}\) вниз; \(16=8t^2/2\), следовательно \(t=\boxed{2\unit{с}}.\)\\
\bottomrule
\end{tabularx}
\vspace{10pt}
\smallhead{Ключевая самопроверка}
Во всех десяти задачах ответ можно получить без запоминания отдельных формул для лифта или нити: достаточно изобразить силы, выбрать ось и аккуратно записать \(\sum F_y=ma_y\). В задаче 10 к найденному ускорению добавляется формула равноускоренного движения.
\end{document}
